\documentclass{article}

\usepackage{amsmath, amssymb, amsthm}
\usepackage[a4paper, margin=1in]{geometry}

\title{Modular Arithmetic Basics}
\author{sid}
\date{\today}

\newtheorem{theorem}{Theorem}
\newtheorem{lemma}{Lemma}
\newtheorem{example}{Example}

\newcommand{\Z}{\mathbb{Z}}
\newcommand{\N}{\mathbb{N}}
\newcommand{\Q}{\mathbb{Q}}
\newcommand{\R}{\mathbb{R}}
\newcommand{\eps}{\varepsilon}

\newcommand{\bld}[1]{\textbf{#1}}
\newcommand{\itl}[1]{\textit{#1}}

\newcommand{\poly}{\mathrm{poly}}
\newcommand{\polylog}{\mathrm{polylog}}
\newcommand{\bigO}{\mathcal{O}}

\begin{document}

\maketitle

\section{What Does ``mod'' Even Mean?}

\subsection*{Definition}

\begin{align*}
    a \equiv b \pmod{n}
    \quad \iff \quad
    n \mid (a - b)
\end{align*}

\textbf{Interpretation:} $a$ and $b$ leave the same remainder when divided by $n$.

\begin{example}
\[
27 \equiv 12 \pmod{5}
\quad \implies \quad
5 \mid (27 - 12)
\]
\end{example}

\subsection*{Basic Properties}

If $a \equiv b \pmod{n}$ and $c \equiv d \pmod{n}$, then:

\begin{align*}
    a + c &\equiv b + d \pmod{n} \\
    a - c &\equiv b - d \pmod{n} \\
    ac    &\equiv bd \pmod{n}
\end{align*}

\textbf{Interpretation:}  
You can do arithmetic normally, then take modulo.

\subsection*{Reduction Rule}

\begin{align*}
    a \equiv a \bmod n
\end{align*}

This means every number can be reduced to its remainder modulo $n$.

\begin{example}
\[
38 \equiv 3 \pmod{5}
\]
\end{example}

\subsection*{Working Modulo $n$}

All computations happen inside:
\[
\mathbb{Z}_n = \{0, 1, 2, \dots, n-1\}
\]

\begin{example}
\[
7 + 9 \equiv 16 \equiv 1 \pmod{5}
\]
\end{example}

\subsection*{Divisibility in Modular Form}

\begin{align*}
    a \equiv 0 \pmod{n}
    \quad \iff \quad
    n \mid a
\end{align*}

\begin{example}
\[
20 \equiv 0 \pmod{5}
\]
\end{example}

\subsection*{Cancellation Rule (Important)}

If:
\[
ac \equiv bc \pmod{n}
\]

Then you can cancel $c$ only if:
\[
\gcd(c, n) = 1
\]

\subsection*{Key Insight}

Modular arithmetic allows us to replace large numbers with smaller representatives while preserving arithmetic structure.


\end{document}